%% =========================================================
%% Chapter: A compactly supported spectral filter
%% Source: OpenAI, "A two-dimensional area law from a global spectral gap"
%%         (September 24, 2026), Lemma 4.2, 03-quasilocal.tex lines 135-205
%% =========================================================

\chapter{A Compactly Supported Spectral Filter}
\label{ch:spectral_filter}

Quasi-local spectral filters remove the matrix elements of an operator
between the ground space and the excited states of a gapped Hamiltonian
exactly, while their time kernel decays fast enough to be combined with
Lieb--Robinson bounds. This chapter constructs the filter
of~\cite[Lemma~4.2]{OpenAI2026AreaLaw}: a smooth cutoff with compact
frequency support whose inverse Fourier transform decays as a stretched
exponential with exponent $p/(p+1)$, for every integer $p\ge1$.

Throughout, $p\ge1$ is an integer, $\alpha=p/(p+1)$, and $\delta>0$.

\section{The profile and its Cauchy estimate}
\label{sec:spectral_filter_profile}

\begin{lemma}[Real part on a small disk]
    \label{lem:spectral_filter_re_bound}
    \lean{SpectralFilter.half_le_re_inv_pow,
        SpectralFilter.re_inv_pow_one_sub_sq_ge}
    \leanok
    If $a\in\C$ satisfies $|a-1|\le 1/(4p)$, then
    $\operatorname{Re}a^{-p}\ge 1/2$. Consequently, for real
    $u\in(-1,1)$, $w=1-u^2$ and $|z-u|\le w/(12p)$, one has $1-z^2\ne0$
    and $\operatorname{Re}(1-z^2)^{-p}\ge w^{-p}/2$.
\end{lemma}
\begin{proof}\leanok
    Put $\varepsilon=1/(4p)$. Bernoulli's inequality gives
    $|a|^p\ge(1-\varepsilon)^p\ge 3/4$ and $(1+\varepsilon)^p\le 4/3$, and
    the geometric sum gives $|a^p-1|\le p\,\varepsilon\,(1+\varepsilon)^p\le1/3$.
    Hence $|a^{-p}-1|\le4/9$. For the second claim write
    $1-z^2=w\,a$ with $a-1=(u-z)(u+z)/w$, so that
    $|a-1|\le 3|z-u|/w\le1/(4p)$.
\end{proof}

\begin{theorem}[Cauchy estimate for the profile]
    \label{thm:spectral_filter_profile_cauchy}
    \lean{SpectralFilter.norm_iteratedDeriv_profile_le_cauchy,
        SpectralFilter.norm_iteratedDeriv_profile_le}
    \leanok
    \uses{lem:spectral_filter_re_bound}
    Let $g(z)=e\exp[-(1-z^2)^{-p}]$, analytic away from $z=\pm1$.
    For $u\in(-1,1)$, $w=1-u^2$ and every $n\ge0$,
    \[
        |g^{(n)}(u)|\le n!\,e\,e^{-w^{-p}/2}\,(12p)^n w^{-n},
    \]
    and for every integer $k$ with $n\le pk$,
    \[
        |g^{(n)}(u)|\le n!\,e\,(12p)^n\,2^k\,k!\;w^{pk-n}.
    \]
\end{theorem}
\begin{proof}\leanok
    The first bound is Cauchy's estimate on the circle of radius $w/(12p)$
    about $u$, where $|g|\le e\,e^{-w^{-p}/2}$ by
    Lemma~\ref{lem:spectral_filter_re_bound}. For the second, put
    $s=w^{-p}$, so that $w^{-n}=w^{pk-n}s^k$, and use
    $s^ke^{-s/2}\le 2^k k!$, which follows from the exponential series.
\end{proof}

\section{The cutoff}
\label{sec:spectral_filter_cutoff}

\begin{definition}[Spectral cutoff]
    \label{def:spectral_cutoff}
    \lean{SpectralFilter.spectralCutoff}
    \leanok
    Define $\chi(\omega)=e\exp[-(1-(\omega/\delta)^2)^{-p}]$ for
    $|\omega|<\delta$ and $\chi(\omega)=0$ for $|\omega|\ge\delta$.
\end{definition}

\begin{theorem}[Smoothness and derivative bounds of the cutoff]
    \label{thm:spectral_cutoff_smooth}
    \lean{SpectralFilter.contDiff_spectralCutoff,
        SpectralFilter.hasCompactSupport_spectralCutoff,
        SpectralFilter.spectralCutoff_zero, SpectralFilter.spectralCutoff_neg,
        SpectralFilter.spectralCutoff_nonneg, SpectralFilter.spectralCutoff_le_one,
        SpectralFilter.abs_iteratedDeriv_spectralCutoff_le}
    \leanok
    \uses{def:spectral_cutoff, thm:spectral_filter_profile_cauchy}
    The cutoff $\chi$ is infinitely differentiable on $\R$, even, supported
    in $[-\delta,\delta]$, takes values in $[0,1]$, and $\chi(0)=1$. For all
    integers $n,k\ge0$ with $n\le pk$ and all $\omega$,
    \[
        |\chi^{(n)}(\omega)|\le\delta^{-n}\,n!\,e\,(12p)^n\,2^k\,k!.
    \]
\end{theorem}
\begin{proof}\leanok
    Let $D_n$ be the $n$-th derivative of $g$ on $(-1,1)$, extended by zero.
    By Theorem~\ref{thm:spectral_filter_profile_cauchy} with $k=n+1$,
    $|D_n(u)|\le C_n|1-u^2|$, so each $D_n$ is continuous at $\pm1$. On
    $(-1,1)$ the derivative of $D_n$ is $D_{n+1}$, outside $[-1,1]$ both
    vanish, and at $\pm1$ a function continuous at a point whose derivative
    extends continuously there is differentiable there with that derivative.
    Hence $D_0$ is infinitely differentiable with $n$-th derivative $D_n$,
    and $\chi(\omega)=D_0(\omega/\delta)$. The bound is the
    Gevrey estimate $\|\chi^{(m)}\|_\infty\le C_1^{m+1}(m!)^{1+1/p}$
    of~\cite[Lemma~4.2]{OpenAI2026AreaLaw}, written with integer exponents.
\end{proof}

\section{The time kernel}
\label{sec:spectral_filter_kernel}

\begin{definition}[Time kernel]
    \label{def:spectral_kernel}
    \lean{SpectralFilter.spectralKernel}
    \leanok
    \uses{def:spectral_cutoff}
    Define $f(t)=(2\pi)^{-1}\int_{\R}\chi(\omega)\cos(t\omega)\,d\omega$.
\end{definition}

\begin{lemma}[Moments of the kernel]
    \label{lem:spectral_kernel_moments}
    \lean{SpectralFilter.spectralKernel_eq_integral_exp,
        SpectralFilter.spectralKernel_neg,
        SpectralFilter.pow_abs_mul_abs_spectralKernel_le}
    \leanok
    \uses{def:spectral_kernel, thm:spectral_cutoff_smooth}
    The kernel is real and even, $f(t)=(2\pi)^{-1}\int\chi(\omega)
    e^{-it\omega}\,d\omega$, and for $n\le pk$,
    \[
        |t|^n|f(t)|\le\frac{\delta}{\pi}\,\delta^{-n}\,n!\,e\,(12p)^n\,
        2^k\,k!.
    \]
\end{lemma}
\begin{proof}\leanok
    Evenness of $\chi$ removes the sine part of $e^{-it\omega}$.
    Integrating by parts $n$ times, with no boundary terms, gives
    $(it)^n\int\chi e^{-it\omega}=\int\chi^{(n)}e^{-it\omega}$; the
    integrand is supported in $[-\delta,\delta]$ and bounded by
    Theorem~\ref{thm:spectral_cutoff_smooth}.
\end{proof}

\begin{theorem}[Stretched-exponential decay]
    \label{thm:spectral_kernel_decay}
    \lean{SpectralFilter.abs_spectralKernel_le_exp,
        SpectralFilter.exists_abs_spectralKernel_le_exp,
        SpectralFilter.integrable_pow_mul_spectralKernel,
        SpectralFilter.integrable_spectralKernel}
    \leanok
    \uses{lem:spectral_kernel_moments}
    For all real $t$,
    $|f(t)|\le C e^{-c|t|^\alpha}$ with $C=2\delta e/\pi$ and
    $c=\log2\cdot\kappa_p\,\delta^\alpha$, where
    $\kappa_p=(4(12p^2)^p)^{-1/(p+1)}$. Every polynomial moment
    $t\mapsto t^nf(t)$ is integrable.
\end{theorem}
\begin{proof}\leanok
    Put $T=\delta|t|$ and $B_p=2(12p^2)^p$. With $n=pk$,
    Lemma~\ref{lem:spectral_kernel_moments}, $(pk)!\le(pk)^{pk}$ and
    $k!\le k^k$ give $T^{pk}|f(t)|\le(\delta e/\pi)(B_pk^{p+1})^k$. Choosing
    $k=\lfloor\kappa_pT^\alpha\rfloor$ makes $B_pk^{p+1}\le T^p/2$, hence
    $|f(t)|\le(\delta e/\pi)2^{-k}\le Ce^{-c|t|^\alpha}$; the case $k=0$
    covers bounded $t$. The moment bounds give
    $(1+|t|)^N|f(t)|\le A_N$ for every $N$, whence integrability.
\end{proof}

\begin{theorem}[Compactly supported spectral filter]
    \label{thm:quasilocal_filter}
    \lean{SpectralFilter.quasilocal_filter,
        SpectralFilter.integral_spectralKernel_mul_exp,
        SpectralFilter.integral_spectralKernel}
    \leanok
    \uses{thm:spectral_cutoff_smooth, lem:spectral_kernel_moments,
        thm:spectral_kernel_decay}
    \cite[Lemma~4.2]{OpenAI2026AreaLaw}.
    The cutoff $\chi$ is smooth and compactly supported with $\chi(0)=1$; the
    real even function $f(t)=(2\pi)^{-1}\int\chi(\omega)e^{-it\omega}\,d\omega$
    satisfies
    \[
        \chi(\omega)=\int_{\R} f(t)e^{it\omega}\,dt,\qquad
        |f(t)|\le Ce^{-c|t|^\alpha},
    \]
    with $C,c>0$ depending only on $p$ and $\delta$. In particular
    $\int_{\R} f=1$.
\end{theorem}
\begin{proof}\leanok
    The Fourier transform of $\chi$ at $\xi$ is $2\pi f(2\pi\xi)$, which is
    integrable by Theorem~\ref{thm:spectral_kernel_decay}. Fourier inversion
    for a continuous integrable function with integrable transform, followed
    by the substitution $t=2\pi\xi$, gives the first identity; $\omega=0$
    gives $\int f=\chi(0)=1$.
\end{proof}

\section{Spectral filtering of a finite matrix}
\label{sec:spectral_filter_matrix}

\begin{definition}[Filtered operator]
    \label{def:filter_integral}
    \lean{SpectralFilter.filterIntegral}
    \leanok
    \uses{def:spectral_kernel}
    For complex square matrices $H$ and $A$, set
    $\mathcal F_H(A)=\int_{\R} f(t)\,e^{itH}Ae^{-itH}\,dt$, a Bochner integral
    for the operator norm.
\end{definition}

\begin{theorem}[Basic properties of the filtered operator]
    \label{thm:filter_integral_basic}
    \lean{SpectralFilter.integrable_filterIntegrand,
        SpectralFilter.norm_filterIntegral_le,
        SpectralFilter.filterIntegral_add, SpectralFilter.filterIntegral_smul,
        SpectralFilter.filterIntegral_finset_sum,
        SpectralFilter.filterIntegral_self,
        SpectralFilter.filterIntegral_conjTranspose,
        SpectralFilter.isHermitian_filterIntegral}
    \leanok
    \uses{def:filter_integral, thm:quasilocal_filter}
    Let $H$ be Hermitian. The integral defining $\mathcal F_H(A)$ converges in
    norm, $\mathcal F_H$ is linear, $\|\mathcal F_H(A)\|\le\|f\|_1\|A\|$,
    $\mathcal F_H(A)^\dagger=\mathcal F_H(A^\dagger)$, and
    $\mathcal F_H(H)=H$.
\end{theorem}
\begin{proof}\leanok
    Since $e^{itH}$ is unitary, $\|e^{itH}Ae^{-itH}\|=\|A\|$, and
    $f$ is integrable by Theorem~\ref{thm:spectral_kernel_decay}. The adjoint
    passes through the integral because $f$ is real. Finally $e^{itH}$
    commutes with $H$, so $\mathcal F_H(H)=(\int f)\,H=H$.
\end{proof}

\begin{theorem}[Spectral multiplier]
    \label{thm:filter_integral_multiplier}
    \lean{SpectralFilter.dotProduct_filterIntegral_mulVec,
        SpectralFilter.star_eigenvectorUnitary_mul_filterIntegral_mul_apply,
        SpectralFilter.filterIntegral_mulVec_of_gap}
    \leanok
    \uses{def:filter_integral, thm:quasilocal_filter}
    Let $H$ be Hermitian with $Hv=Ev$ and $Hw=E'w$. Then
    $\langle v,\mathcal F_H(A)w\rangle=\chi(E-E')\langle v,Aw\rangle$. In an
    orthonormal eigenbasis $(u_k)$ of $H$ with eigenvalues $\lambda_k$, the
    entries of $\mathcal F_H(A)$ are those of $A$ multiplied by
    $\chi(\lambda_k-\lambda_l)$. If moreover every eigenvalue $\lambda_k$
    equals $E'$ or satisfies $|\lambda_k-E'|\ge\delta$, then
    $\mathcal F_H(A)w=\sum_{\lambda_k=E'}\langle u_k,Aw\rangle u_k$, the
    projection of $Aw$ onto the $E'$-eigenspace.
\end{theorem}
\begin{proof}\leanok
    The exponential series gives $e^{-itH}w=e^{-itE'}w$ and
    $\langle v,e^{itH}y\rangle=e^{itE}\langle v,y\rangle$, so the integrand
    has matrix element $f(t)e^{it(E-E')}\langle v,Aw\rangle$, and Fourier
    inversion (Theorem~\ref{thm:quasilocal_filter}) evaluates the integral.
    For the last claim, $\chi(0)=1$ and $\chi$ vanishes outside
    $(-\delta,\delta)$, so the coefficients of $\mathcal F_H(A)w$ in the
    eigenbasis are $\langle u_k,Aw\rangle$ when $\lambda_k=E'$ and zero
    otherwise.
\end{proof}

\section{Positive replacement of a gapped Hamiltonian}
\label{sec:positive_replacement}

Throughout, $H=\sum_i h_i$ is a finite sum of Hermitian matrices with
$\|h_i\|\le J$, $\Omega$ is a unit vector with $H\Omega=E_0\Omega$, and
$H-E_0I\ge\Delta(I-P_\Omega)$ with $\Delta>0$, where
$P_\Omega=|\Omega\rangle\langle\Omega|$.

\begin{lemma}[Filtering below the gap]
    \label{lem:filter_below_gap}
    \lean{SpectralFilter.filterIntegral_mulVec_of_posSemidef_gap}
    \leanok
    \uses{thm:filter_integral_multiplier}
    If $0<\delta\le\Delta$, then $\mathcal F_H(A)\Omega=\langle\Omega,A\Omega\rangle\Omega$
    for every matrix $A$.
\end{lemma}
\begin{proof}\leanok
    Test against an orthonormal eigenbasis $(u_k)$ of $H$. If
    $\lambda_k=E_0$, the gap inequality at $u_k$ gives
    $\Delta(1-|\langle\Omega,u_k\rangle|^2)\le0$, hence
    $\|u_k-\langle\Omega,u_k\rangle\Omega\|^2=1-|\langle\Omega,u_k\rangle|^2=0$,
    and $\chi(0)=1$. If $\lambda_k\ne E_0$, then $u_k\perp\Omega$, so the gap
    inequality gives $\lambda_k-E_0\ge\Delta\ge\delta$ and $\chi$ vanishes.
\end{proof}

\begin{theorem}[Centered filtered terms]
    \label{thm:centered_filter}
    \lean{SpectralFilter.centeredFilter, SpectralFilter.isHermitian_centeredFilter,
        SpectralFilter.centeredFilter_mulVec_eq_zero,
        SpectralFilter.sum_centeredFilter, SpectralFilter.norm_centeredFilter_le}
    \leanok
    \uses{lem:filter_below_gap, thm:filter_integral_basic}
    Let $0<\delta\le\Delta$ and
    $M_i=\mathcal F_H(h_i)-\langle\Omega,h_i\Omega\rangle I$. Then $M_i$ is
    Hermitian, $M_i\Omega=0$, $\sum_iM_i=H-E_0I$, and
    $\|M_i\|\le\|h_i\|(\|f\|_1+1)$.
\end{theorem}
\begin{proof}\leanok
    Lemma~\ref{lem:filter_below_gap} gives $M_i\Omega=0$. Linearity and
    $\mathcal F_H(H)=H$ give the sum, since
    $\sum_i\langle\Omega,h_i\Omega\rangle=\langle\Omega,H\Omega\rangle=E_0$.
    The norm bound combines $\|\mathcal F_H(h_i)\|\le\|f\|_1\|h_i\|$ with
    $|\langle\Omega,h_i\Omega\rangle|\le\|h_i\|$.
\end{proof}

\begin{lemma}[Square-root estimate]
    \label{lem:sqrt_holder}
    \lean{CFC.norm_sqrt_sub_sqrt_le, CFC.le_abs_of_isSelfAdjoint,
        CFC.norm_abs_sub_abs_le}
    \leanok
    In a unital C*-algebra, $\|a^{1/2}-b^{1/2}\|\le\|a-b\|^{1/2}$ for
    $a,b\ge0$. For self-adjoint $a,b$ one has $a\le|a|$ and
    $\||a|-|b|\|\le((\|a\|+\|b\|)\|a-b\|)^{1/2}$.
\end{lemma}
\begin{proof}\leanok
    With $\epsilon=\|a-b\|$, operator monotonicity of the square root gives
    $a^{1/2}\le(b+\epsilon)^{1/2}\le b^{1/2}+\epsilon^{1/2}$, the last step
    because $b+\epsilon\le(b^{1/2}+\epsilon^{1/2})^2$. Exchanging $a$ and $b$
    bounds the norm. For the absolute values apply this to $a^2,b^2$ and use
    $a^2-b^2=a(a-b)+(a-b)b$.
\end{proof}

\begin{theorem}[Positive replacement, spectral part]
    \label{thm:positive_replacement}
    \lean{SpectralFilter.positiveConstraint, SpectralFilter.positiveNormalization,
        SpectralFilter.positiveConstraint_nonneg,
        SpectralFilter.positiveConstraint_le_one,
        SpectralFilter.positiveConstraint_mulVec_eq_zero,
        SpectralFilter.smul_le_positiveConstraint,
        SpectralFilter.positive_replacement}
    \leanok
    \uses{thm:centered_filter, lem:sqrt_holder}
    Take $\delta=\Delta/2$, $c_*=\max\{1,J(\|f\|_1+1)\}$, and
    $k_i=|M_i|/c_*$. Then $0\le k_i\le I$, $k_i\Omega=0$, and
    \[
        \sum_ik_i\ge c_*^{-1}(H-E_0I)\ge\frac{\Delta}{c_*}(I-P_\Omega).
    \]
\end{theorem}
\begin{proof}\leanok
    Theorem~\ref{thm:centered_filter} gives $\|M_i\|\le c_*$, hence
    $\|k_i\|\le1$. Since $|M_i|^2=M_i^2$, the kernel of $M_i$ lies in that of
    $|M_i|$, so $k_i\Omega=0$. Finally $M_i\le|M_i|$ and $\sum_iM_i=H-E_0I$.
\end{proof}

\section{Square-root channels}
\label{sec:root_channels}

\begin{theorem}[Square roots and root channels]
    \label{thm:root_channel}
    \lean{Matrix.rootChannel, Matrix.norm_sum_conjTranspose_mul_mul_le,
        Matrix.l2_opNorm_kronecker_one_le, Matrix.rootChannel_one,
        Matrix.rootChannel_nonneg, Matrix.rootChannel_kronecker_nonneg,
        Matrix.norm_rootChannel_le,
        Matrix.norm_rootChannel_kronecker_le, Matrix.norm_rootChannel_sub_le,
        Matrix.norm_rootChannel_kronecker_sub_le,
        Matrix.sqrt_one_sub_mul_add_sqrt_mul, Matrix.norm_sqrt_sub_le,
        Matrix.norm_sqrt_one_sub_sub_le, Matrix.sqrt_mulVec_eq_zero,
        Matrix.sqrt_one_sub_mulVec_eq_self,
        Matrix.rootChannel_sqrt_eq_self_of_commute}
    \leanok
    \uses{lem:sqrt_holder}
    Let $0\le k,k'\le I$, $G=(I-k)^{1/2}$, $K=k^{1/2}$, and define $G',K'$
    from $k'$. Then $G^2+K^2=I$, and
    $\|G-G'\|,\|K-K'\|\le\|k-k'\|^{1/2}$; if $k\Omega=0$, then $K\Omega=0$
    and $G\Omega=\Omega$. The map $\mathcal E(B)=GBG+KBK$ is unital and
    completely positive, and for every finite auxiliary system the extension
    $\mathcal E\otimes\mathrm{id}$ is a contraction in operator norm.
    Moreover, uniformly in the auxiliary dimension,
    \[
        \|(\mathcal E-\mathcal E')\otimes\mathrm{id}\|
        \le(\|G\|+\|G'\|)\|G-G'\|+(\|K\|+\|K'\|)\|K-K'\|,
    \]
    and $\mathcal E$ fixes every matrix commuting with $k$.
\end{theorem}
\begin{proof}\leanok
    The map $v\mapsto Gv\oplus Kv$ is an isometry, and $\mathcal E(B)$ is the
    compression of $B\oplus B$ by it; tensoring with the identity preserves
    $G^2+K^2=I$. The difference bound follows from
    $ABA-A'BA'=(A-A')BA+A'B(A-A')$ and $\|A\otimes I\|\le\|A\|$. The root
    tails are Lemma~\ref{lem:sqrt_holder}. A matrix commuting with $k$
    commutes with both square roots, so its image is itself times
    $G^2+K^2=I$.
\end{proof}

\section{Locality of the positive constraints}
\label{sec:positive_locality}

Let $L$ be a linear map on matrices, playing the role of a conditional
expectation onto operators supported in a ball of radius $l$, and write
$\tau_t(A)=e^{itH}Ae^{-itH}$ and $\alpha=p/(p+1)$.

\begin{lemma}[Kernel tail]
    \label{lem:spectral_kernel_tail}
    \lean{SpectralFilter.abs_spectralKernel_le_tail,
        SpectralFilter.integrable_inv_one_add_abs_sq}
    \leanok
    \uses{thm:spectral_kernel_decay}
    There is $K_0>0$ with
    $|f(t)|\le K_0(1+|t|)^{-2}e^{-(c/2)T^\alpha}$ whenever $0\le T\le|t|$,
    where $c$ is the decay rate of the kernel.
\end{lemma}
\begin{proof}\leanok
    Multiply the bounds $|f(t)|\le A(1+|t|)^{-4}$ and
    $|f(t)|\le Ke^{-c|t|^\alpha}$ and take square roots.
\end{proof}

\begin{theorem}[Localization of filtered operators]
    \label{thm:filter_localization}
    \lean{SpectralFilter.exists_norm_filterIntegral_sub_map_le',
        SpectralFilter.exists_norm_filterIntegral_sub_map_le}
    \leanok
    \uses{lem:spectral_kernel_tail, def:filter_integral}
    Let $C\ge0$ and $v,c>0$. There are $C',c'>0$, depending only on
    $p,\delta,C,v,c$, such that
    $\|\tau_t(A)-L\tau_t(A)\|\le\|A\|\min\{2,Ce^{v|t|-cl}\}$ for all $t$
    implies $\|\mathcal F_H(A)-L\mathcal F_H(A)\|\le C'e^{-c'l^\alpha}\|A\|$.
\end{theorem}
\begin{proof}\leanok
    The linear map $L$ passes through the integral. For $|t|\le cl/(2v)$ the
    hypothesis gives the factor $Ce^{-cl/2}$; for larger $|t|$ use the bound
    $2$ and Lemma~\ref{lem:spectral_kernel_tail}. Finally
    $e^{-cl/2}\le e^{c/2}e^{-(c/2)l^\alpha}$.
\end{proof}

\begin{theorem}[Tail of the positive constraints]
    \label{thm:positive_constraint_tail}
    \lean{SpectralFilter.norm_abs_sub_map_abs_le,
        SpectralFilter.exists_norm_positiveConstraint_sub_map_le}
    \leanok
    \uses{thm:filter_localization, thm:positive_replacement, lem:sqrt_holder}
    If $L$ is a contraction, $M$ and $LM$ are Hermitian, and $L$ fixes $|LM|$,
    then $\||M|-L|M|\|\le2(2\|M\|\|M-LM\|)^{1/2}$. Consequently, if $L$ is
    moreover unital and preserves Hermitian matrices, and $h$ satisfies the
    hypothesis of Theorem~\ref{thm:filter_localization} with $\|h\|\le J$,
    then the positive constraint $k=|M|/c_*$ built from $h$ satisfies
    $\|k-Lk\|\le C_3e^{-c_3l^\alpha}$, with constants depending only on
    $p,\Delta,J,C,v,c$.
\end{theorem}
\begin{proof}\leanok
    Write $|M|-L|M|=(|M|-|LM|)+L(|LM|-|M|)$ and apply
    Lemma~\ref{lem:sqrt_holder}. Unitality gives $M-LM=\mathcal F_H(h)-L\mathcal F_H(h)$.
\end{proof}

\section{Stability of a global gap}
\label{sec:gap_perturbation}

\begin{lemma}[Trace distance of pure states]
    \label{lem:pure_trace_distance}
    \lean{Matrix.abs_eq_inv_sqrt_smul_mul_self,
        Matrix.sub_mul_sub_mul_sub_of_idempotent, Matrix.traceDistance_vecMulVec_eq}
    \leanok
    For unit vectors $\psi,\phi$ the trace distance of
    $|\psi\rangle\langle\psi|$ and $|\phi\rangle\langle\phi|$ is
    $(1-|\langle\psi,\phi\rangle|^2)^{1/2}$.
\end{lemma}
\begin{proof}\leanok
    With $X=P-Q$ and $s=|\langle\psi,\phi\rangle|^2$ one has $PQP=sP$ and
    $QPQ=sQ$, hence $X^3=(1-s)X$; so $|X|=X^2/(1-s)^{1/2}$ and
    $\operatorname{tr}X^2=2(1-s)$.
\end{proof}

\begin{theorem}[Stability of a global gap]
    \label{thm:gap_perturbation}
    \lean{exists_unit_eigenvector_re_inner_le, re_inner_smul_add_eigenvector,
        exists_gap_of_norm_sub_le, Matrix.exists_posSemidef_gap_of_norm_sub_le}
    \leanok
    \uses{lem:pure_trace_distance}
    Let $H_F\ge g(I-P_\Omega)$ with $H_F\Omega=0$, $\|\Omega\|=1$, $g>0$, and
    let $H'\ge0$ with $\|H'-H_F\|\le\epsilon\le g/4$. Then $H'$ has a unit
    ground vector $\Omega_0$ with energy $e_0\in[0,\epsilon]$ and
    $H'-e_0I\ge(g/2)(I-P_{\Omega_0})$; after a choice of phase
    $\|\Omega_0-\Omega\|\le2(\epsilon/g)^{1/2}$, and the trace distance of the
    two pure states is at most $(2\epsilon/g)^{1/2}$.
\end{theorem}
\begin{proof}\leanok
    A minimizer of the quadratic form on the unit sphere is an eigenvector.
    For $\psi\perp\Omega_0$, a combination of $\psi$ and $\Omega_0$ orthogonal
    to $\Omega$ has energy at least $(g-\epsilon)$ times its squared norm,
    while the $\Omega_0$ component carries energy $e_0\le g-\epsilon$; hence
    $\langle\psi,H'\psi\rangle\ge(g-\epsilon)\|\psi\|^2$. Testing on $\Omega$
    gives $e_0\le\epsilon$, and the gap of $H_F$ at $\Omega_0$ gives
    $g(1-|\langle\Omega,\Omega_0\rangle|^2)\le2\epsilon$.
\end{proof}
